\begin{frame}{4단계로 뜯어보기}
  \begin{enumerate}
    \item \kb{$QK^T$} — 모든 단어 쌍의 Query-Key 내적을 \textbf{한
      번에} 계산한 행렬. 각 행 = 한 단어가 다른 모든 단어와
      얼마나 관련 있는지에 대한 \textbf{원점수}.
    \item \kb{$\sqrt{d_k}$ 로 나누기} (Key 벡터 차원의 제곱근) —
      내적값이 차원이 커질수록 커지는 경향이 있어, softmax가
      극단적인 값(0 또는 1에 가까운)만 내뱉게 되는 것을 막는
      \textbf{스케일 조정}.
    \item \kb{softmax} — 각 행을 \textbf{확률분포}(합이 1)로
      바꿈. ``이 단어가 다른 단어들에게 나눠 주는 주의의 비율''.
    \item \kb{$V$ 가중합} — 그 확률로 $V$(각 단어가 실제로
      전달할 내용)의 가중합을 냄. 결과 = \textbf{``관련 단어들의
      내용을 관련도만큼 섞어 담은''} 새로운 벡터.
  \end{enumerate}
\end{frame}
